\section{利用三角板画心脏线}
\markboth{29. 利用三角板画心脏线}{29. 利用三角板画心脏线}

只要先画一个辅助圆，就能利用三角板画出一条心脏线来。

\noindent
\begin{minipage}[c]{0.42\textwidth}
  \centering
  \begin{tikzpicture}[scale=0.80, >=Stealth, line cap=round, line join=round]
    % 圆 C 半径 2r = 1.0, O=(0,0), C=(1.0, 0), theta = 52 deg
    \def\rr{1.0}
    \def\th{52}
    \coordinate (O) at (0,0);
    \coordinate (C) at (\rr, 0);
    \coordinate (B) at ({\rr*cos(\th)*cos(\th)}, {\rr*cos(\th)*sin(\th)});
    \coordinate (M) at ({\rr*(1+cos(\th))*cos(\th)}, {\rr*(1+cos(\th))*sin(\th)});
    \coordinate (N) at ({\rr*(1+cos(\th))}, {\rr*sin(\th)});

    % 极轴
    \draw[->] (O) -- (2.55, 0);

    % 辅助圆 C 与心脏线
    \draw (C) circle (\rr);
    \draw[very thick, curveblue, domain=0:360, samples=200, variable=\t]
      plot ({\rr*(1+cos(\t))*cos(\t)}, {\rr*(1+cos(\t))*sin(\t)});

    % 矩形 CBMN 与线段 OM
    \draw (O) -- (M);
    \draw (C) -- (B);
    \draw (M) -- (N);
    \draw (C) -- (N);

    % 直角记号 (直角顶点 M 与垂足 B)
    \draw ($(M)+({\th-180}:0.14)$) -- ($(M)+({\th-180}:0.14)+({\th-90}:0.14)$) -- ($(M)+({\th-90}:0.14)$);
    \draw ($(B)+(\th:0.12)$) -- ($(B)+(\th:0.12)+({\th-90}:0.12)$) -- ($(B)+({\th-90}:0.12)$);

    % 关键点与标签
    \fill (O) circle (1.3pt) node[below right=1pt, font=\footnotesize] {$O$};
    \fill (C) circle (1.3pt) node[below, font=\footnotesize] {$C$};
    \fill (B) circle (1.3pt) node[above left=-1pt, font=\footnotesize] {$B$};
    \fill (M) circle (1.3pt) node[above, font=\footnotesize] {$M$};
    \fill[white] (N) circle (1.4pt); \draw (N) circle (1.4pt) node[right, font=\footnotesize] {$N$};
  \end{tikzpicture}
  \captionof*{figure}{图\ 29.1}
\end{minipage}
\hfill
\begin{minipage}[c]{0.55\textwidth}
  \setlength{\parindent}{2em}\indent
  如图 29.1，设辅助圆为圆 $C$。在圆周上任取一点 $O$。让三角板的一条直角边通过 $O$ 点，另一条直角边与圆 $C$ 相切，这时三角板的直角顶点 $M$ 就是心脏线的一个点。继续这样作出心脏线上的若干个点，最后只要用光滑曲线把这些点连起来就行了（但 $O$ 点处有一个尖角）。
\end{minipage}

为了证明上述画法的正确性，我们取 $O$ 为极点，$OC$ 为极轴，并设圆 $C$ 的半径为 $2r$。设三角板的一条直角边沿 $MN$ 与圆 $C$ 相切，切点是 $N$。作 $CB\perp OM$，垂足为 $B$（图 29.1），则 $CBMN$ 是矩形。因而
\begin{align*}
  \rho = OM &= OB + BM \\
            &= OC\cos\theta + CN \\
            &= 2r\cos\theta + 2r,
\end{align*}
即
\[
  \rho = 2r(1+\cos\theta),
\]
这正是心脏线的极坐标方程（见第 28 节）。

利用第 10 节中介绍的垂足曲线概念，上面所证明的事实可以叙述成：心脏线是一个圆关于圆周上一点的垂足曲线。
